% !TEX root = ../../../main.tex
\section{Permutations}
\label{sec:linear-algebra-i:permutations}

% Sources: October 1 PDF, pp. 20--24; IMG_9208--IMG_9210.
\subsection{Permutations and composition}

\begin{definition}
  A permutation of \(\{1,\ldots,n\}\) is a bijective map
  % Author-requested combinatorial context, October 6, 2026.
  \marginnote{%
    \textbf{Permutations in combinatorics.}
    A permutation is an ordered arrangement of distinct objects, without
    repetition. The list \((\sigma(1),\ldots,\sigma(n))\) identifies a full
    arrangement with a bijection in \(S_n\). There are \(n\) choices for the
    first entry, \(n-1\) for the next, and so on, giving \(n!\) arrangements.
    \par
    For \(0\le k\le n\), the binomial coefficient
    \[
      \binom{n}{k}=\frac{n!}{k!(n-k)!}
    \]
    counts unordered \(k\)-element subsets. Each subset has \(k!\) orderings,
    so the number of ordered selections of \(k\) distinct objects is
    \[
      k!\binom{n}{k}=\frac{n!}{(n-k)!},
    \]
    with \(0!=1\). These same binomial coefficients appear in \((x+y)^n\).
    Choosing \(k\) of the \(n\) factors to contribute \(y\) gives the term
    \(\binom{n}{k}x^{n-k}y^k\).
  }
  \[
    \sigma:\{1,\ldots,n\}\longrightarrow\{1,\ldots,n\}.
  \]
  Write \(S_n\) for the set of all such permutations.
\end{definition}

Permutations reorder \(n\) objects. Under composition, \(S_n\) is a group.
The identity permutation \(\iota\) belongs to \(S_n\); the composition
of two permutations is a permutation; and every \(\sigma\in S_n\) has
an inverse satisfying
\[
  \sigma\circ\sigma^{-1}=\sigma^{-1}\circ\sigma=\iota.
\]
Composition is associative.

\begin{definition}
  A transposition switches two elements and leaves all the others unchanged.
  For \(i\ne j\), the transposition \(\tau_{ij}\in S_n\) is
  \[
    \tau_{ij}(k)=
    \begin{cases}
      j & \text{if \(k=i\)},\\
      i & \text{if \(k=j\)},\\
      k & \text{otherwise}.
    \end{cases}
  \]
\end{definition}

\begin{example}
  The transposition \(\tau_{24}\in S_4\) satisfies
  \[
    \tau_{24}(1)=1,\quad \tau_{24}(2)=4,\quad
    \tau_{24}(3)=3,\quad \tau_{24}(4)=2.
  \]
\end{example}

Every permutation can be written as a composition of transpositions,
\[
  \sigma=\tau_k\circ\cdots\circ\tau_1.
\]
The identity is the composition of zero transpositions.
% PDF pp. 20--21 and IMG_9209 assert decomposition without a proof.

\begin{example}
  The permutation \(\sigma\in S_3\) given by
  \[
    \sigma(1)=2,\qquad \sigma(2)=3,\qquad \sigma(3)=1
  \]
  is \(\sigma=\tau_{13}\circ\tau_{12}\). Here
  \[
    \begin{aligned}
      \tau_{12}(1)&=2,& \tau_{12}(2)&=1,& \tau_{12}(3)&=3,\\
      \tau_{13}(1)&=3,& \tau_{13}(2)&=2,& \tau_{13}(3)&=1.
    \end{aligned}
  \]
\end{example}

\subsection{Parity and the sign function}

\begin{proposition}
  \label{thm:linear-algebra-i:permutation-parity}
  A permutation cannot be both a composition of an even number of
  transpositions and a composition of an odd number of transpositions.
\end{proposition}
% PDF p. 23 and IMG_9209 state this result without supplying a proof.

Consequently, the following definition does not depend on the chosen
transposition decomposition.

\begin{definition}
  The sign, or parity, of a permutation is the function
  \(\varepsilon:S_n\to\{-1,1\}\). For a decomposition
  \(\sigma=\tau_k\circ\cdots\circ\tau_1\) into transpositions,
  \[
    \varepsilon(\sigma)=
    \begin{cases}
      1 & \text{if \(k\) is even},\\
      -1 & \text{if \(k\) is odd}.
    \end{cases}
  \]
\end{definition}

The sign function has the properties
\[
  \begin{aligned}
    \varepsilon(\iota)&=1,\\
    \varepsilon(\tau)&=-1 &&\text{for every transposition \(\tau\)},\\
    \varepsilon(\sigma_2\circ\sigma_1)
      &=\varepsilon(\sigma_2)\varepsilon(\sigma_1),\\
    \varepsilon(\sigma^{-1})&=\varepsilon(\sigma).
  \end{aligned}
\]
For the last property, reversing the order of the same transpositions gives
the inverse,
\[
  \sigma=\tau_1\circ\cdots\circ\tau_N
  \quad\Longrightarrow\quad
  \sigma^{-1}=\tau_N\circ\cdots\circ\tau_1.
\]
The sign function is uniquely determined by assigning sign \(-1\) to
every transposition and requiring multiplicativity under composition.
Its existence depends on \autoref{thm:linear-algebra-i:permutation-parity}.

\subsection{Extending the sign to all maps}

Let \(T_n\) be the set of all maps
\[
  \phi:\{1,\ldots,n\}\longrightarrow\{1,\ldots,n\}.
\]
Then \(S_n\subseteq T_n\). Extend the sign function to
\(\varepsilon:T_n\to\{-1,0,1\}\) by keeping its existing values on
\(S_n\) and setting
\[
  \varepsilon(\phi)=0\qquad\text{if \(\phi\notin S_n\)}.
\]
